% ===========================================================================
% Section 1: Introduction (~1.5 pages)
% ===========================================================================

% --- Problem statement ---
% I/O bottlenecks in HPC: gap between detection and actionable fixes

% --- Existing approaches and limitations ---
Prior work on I/O bottleneck detection includes rule-based
systems~\cite{bez2022drishti, devarajan2025wisio},
ML classifiers~\cite{dong2023aiio}, and LLM-based
approaches~\cite{zheng2025ioagent, lockwood2024ion}.

% TODO: Expand on limitations of each:
% - Rule-based (Drishti, WisIO): Cannot learn, brittle thresholds, no code fixes
% - ML-only (AIIO): Detects patterns but no recommendations
% - LLM-only (IOAgent, ION): No learned detection, hallucination risk, slow

% --- Our approach ---
% ML+LLM hybrid with benchmark-grounded Knowledge Base
% Three-phase architecture: benchmark KB -> ML training -> online inference

% --- Contributions ---
% TODO: 3-4 bullet points
% \begin{itemize}
%   \item A hybrid ML+LLM architecture that combines learned bottleneck
%         detection with retrieval-augmented code-level recommendations.
%   \item A benchmark-grounded Knowledge Base mapping Darshan I/O signatures
%         to verified code causes and validated fixes.
%   \item The first ML-based I/O analysis at facility scale on 1.37M Darshan
%         logs from ALCF Polaris.
%   \item Comprehensive evaluation showing [TODO: key results].
% \end{itemize}
